% NAF 9544, batch 1 — Gallica views 1–20.
% Pass: Opus 5.5 (claude-opus-5-5), 2026-09-22 — first pass, unchecked against the views by a human.
%
% What the volume is. An album bound in 1899 (« Autographes de Savants —
% Volume de 36 Feuillets — 27 Novembre 1899 », v. 3): autograph pieces, each
% mounted on a guard or on a mounting leaf, usually preceded by a dealer's
% printed catalogue cutting or a collector's docket. Colour-free greyscale
% scans, landscape, about 8290 × 5779 (4636 × 5794 for v. 1): each view is the
% album lying open. Every piece carries the mark « R. D. 9555 » with a brace,
% and the round stamp « BIBLIOTHÈQUE NATIONALE — R. F. — MSS ».
%
% What the twenty views hold. Nothing in this batch is of the seventeenth
% century: the catalogue's dating « 1601-1700 » is copied verbatim below, but
% the pieces here run from 1765 to the year X of the Republic; the Huygens and
% Mersenne letters the title promises must lie further on.
%
%   v. 1–2    Front board and marbled pastedown with the label « FR / NOUV.
%             ACQ. / 9544 ». Not transcribed.
%   v. 3      Title leaf (« Autographes de Savants », leaf count and date of
%             1899). Given as a note.
%   v. 4      Mounting leaf, with a printed sale-catalogue cutting announcing
%             the d'Alembert piece. Given as a note.
%   v. 5–6    D'ALEMBERT, minute to Malesherbes, undated on the leaf; a later
%             small hand adds « a M. de Malesherbes » and « 1765 ». Written in
%             the third person (« M. D'alembert est pensionnaire surnumeraire
%             depuis 9 ans et demi dans la classe de Geometrie ») and asking
%             that Malesherbes represent to M. de St. Florentin that he pass
%             to a place of pensionnaire mécanicien. Two small leaves: begins
%             at the head of v. 5, runs onto the verso (v. 6 left) and ends on
%             the second leaf (v. 6 right), « … après les travaux \& les
%             sacrifices de M. d'alembert », with a closing stroke. No
%             signature, no salutation.
%   v. 7      Mounting leaf with the Bossut letter folded down over it: only
%             dealers' figures and the name « Bossut ». Given as a note.
%   v. 8–10   BOSSUT to « Monsieur le Comte », Paris, 12 January 177[4?]; the
%             last figure of the year is open and uncertain (4 rather than 6).
%             Docket « Reponse donnée » in another hand. He thanks the count
%             for his election to « votre celebre academie », says he has
%             passed on the count's message to d'Alembert. Begins v. 8, runs
%             v. 9 (both pages), ends v. 10 with « Monsieur le Comte, votre
%             très humble et très obeissant serviteur Bossut ». The count and
%             his academy are not named on the leaves.
%   v. 11     Mounting leaf with a printed catalogue cutting (Bossut to Antide
%             Janvier, 2 novembre an I) and a pencil note « cat. Dufour 30
%             janv. 1889 » (Dufour uncertain). Given as a note.
%   v. 12     Collector's docket (« C. Bossut. Savant Geometre … ») pasted by
%             its top edge over the next letter; pencil « Arrigoni, XIIIe
%             catalogue 1878 ».
%   v. 13     SKIPPED: a second exposure of v. 12, the docket down, nothing new.
%   v. 14–16  BOSSUT to Janvier, « Paris 2 9bre l'an 1. de la Rep. ». The
%             docket is lifted on v. 14. Thanks for a portrait of Lalande;
%             declines to advise on a petition to the Assemblée nationale;
%             advises him to ask the minister for a lodging at the Menus.
%             Begins v. 14, continues v. 15 left, ends v. 15 right, « Recevés
%             les salutations de votre concitoyen Bossut ». Address panel on
%             v. 16, written across the page: « A Monsieur / Monsieur Janvier,
%             horloger mecanicien. aux menus, rue poissonniere et bergere / A
%             Paris ».
%   v. 17     BOSSUT, a receipt to the « Citoyen agasse » for the 65th–67th
%             livraisons of the Encyclopédie méthodique, « A Paris ce 3
%             fructidor, an 10e », signed. Below it, an Italian collector's
%             wrapper listing Bossut, Laplace and Delambre.
%   v. 18–19  BOSSUT, the one-leaf summary (« sommaire ») of his memoir on the
%             equilibrium of vaults, headed « 21 Messidor, an 9me de la Rep. »:
%             three sections, chapters listed. Two reading notes in the left
%             margin of v. 18: « Lû le 21 messidor an 9 à la réserve des
%             calculs » (a modern pencil note attributes it to Laplace) and
%             « Le 21 messidor an 9 le cit. Bossut a lû le memoire dont cette
%             feuille est le sommaire », signed Delambre. The summary begins at
%             the head of v. 18 and ends on its verso, v. 19, with Section III.
%             Collector's note at the head: « écriture de M. l'abbé Bossut ».
%   v. 20     Wrapper « 677 / Sofia Germain (autogr.) » and a dealer's or
%             collector's note « Sous ce pli l'autographe de Sophie Germain,
%             acheté le lundi, 30 juin, (cabinet Dubrunfaut) … A. M. ». The
%             Germain autograph itself begins on v. 21 (next batch).
%
% Hands. D'Alembert's minute: a fast, sloping cursive, long s throughout, « \& »
% for et; he writes « mechanicien », « mechanique », « Geometrie », « vouloit »,
% « seroit », « puisqu'il », « reciproquement », « veterance ». Bossut: a small
% upright hand whose final r is a small open loop that reads as o or w
% (« Monsieur », « honneur », « tenir »), and whose « Comte » reads « Conte »;
% he writes « avés », « m'aviés », « reconnoissance », « remerciments »,
% « datte », « luy », « paroit ». Long s is set s throughout, said here once.
% Abbreviations kept: Mr, M., St, C. and cit. for citoyen, 9bre, Rep.
%
% Numbers on the leaves. Several series, none settled as the library's:
%  - One ink figure at the top right of each recto, a thin clerk's hand with a
%    flagged 1 and 7, distinct from every writer: 1 (v. 4, the first mounting
%    leaf), 2 (v. 5), 3 (v. 6 right), 4 (v. 7 mounting leaf), 5 (v. 8), 6
%    (v. 9 right), 7 (v. 11), 8 (v. 12 docket), 9 (v. 14), 10 (v. 15 right),
%    11 (v. 17 receipt), 12 (v. 17 wrapper), 13 (v. 18), 14 (v. 20 wrapper),
%    15 (v. 20 note). One number per leaf, counting mounts, dockets and
%    wrappers, across pieces by different writers, consistent with the leaf
%    count of the 1899 certificate. It sits where the foliation sits and
%    behaves as a foliation; but in a greyscale scan the figures look inked,
%    and following the house practice (Latin 17859, Latin 11195) no \folio{}
%    is written. A human look would settle it at once.
%  - Large pencil lot numbers, one per piece: 207 (v. 4), 96 (v. 7, ringed),
%    95 (v. 11), 1498 or 4498 (v. 12), 97 (v. 17, ringed), 677 (v. 20).
%  - Small ink numbers from earlier collections: 40 and 41 on the letter to
%    the count (v. 8, 9); 109 (v. 11), 113 (v. 12), 111 (v. 15), 114 (v. 17);
%    17a (v. 11), 17b (v. 12), 18 and 7 (v. 17); 64 (v. 17), 65 (v. 18), 65/1
%    and 65/2 (v. 20); 6 and 16 (v. 7); 20 (v. 11); « CLXXI 2810 » (v. 20);
%    a lone 2 on the address panel (v. 16).
%  - « R. D. 9555 » on every piece; « N. a. franç. 9544. » on v. 4.
%
% The catalogue's dating is copied verbatim; the dates on the leaves (1765 as
% annotated, 177[4], an I, an 9, an 10) are transcribed where they stand.

\documentclass[11pt,a4paper]{article}
% The preamble shared by every transcription of the mathematicians' manuscripts in Gallica.
%
% It rests on one idea: **uncertainty compounds, it does not resolve.** A
% transcription that smooths over a doubtful word has destroyed the only thing
% it was made for — a reader must be able to tell, at every point, what was on
% the page and what was supplied. The apparatus macros below are therefore the
% critical apparatus, not decoration.
%
% The same commands are recognised by `scripts/render.mjs`, which produces the
% site's reading view, and by `scripts/tei.mjs`. A macro added here must be
% added there too, or rendering fails loudly — which is the wanted behaviour.
%
% Comments here are English, like the rest of the repository. The documents
% this preamble sets are French, and that is deliberate: see the skills.

\ProvidesPackage{germain}[2026/09/15 Transcriptions of mathematicians' manuscripts in Gallica]

\usepackage{amsmath,amssymb,amsthm}
\usepackage{mathrsfs}
% Kept from the parent preamble so the renderer's diagram support stays
% available; these manuscripts draw no commutative diagrams, and nothing here uses it.
\usepackage{tikz-cd}
\usepackage[english,french]{babel}
\usepackage{fontspec}

% --- Greek ------------------------------------------------------------------
% The text face stays fontspec's default, Latin Modern, which has no Greek:
% XeTeX drops every glyph it lacks with a line in the log and nothing on the
% page, and the Greek words the manuscripts quote vanished from the PDFs. So
% the Greek and Greek Extended blocks get a XeTeX character class of their own,
% and entering or leaving a run of it switches the family to CMU Serif — the
% same Computer Modern design, with polytonic Greek — and back. Only the
% family changes: a Greek word in \textit or \textbf keeps its shape and series.
% The transitions to and from every other class carry the tokens that class
% has towards an ordinary letter, so French spacing before « ; » or inside
% « » still holds after a Greek word. No grouping, so no brace or \endgroup
% can come out unbalanced; maths is left alone, since XeTeX inserts nothing
% there. Kept identical in `exercices.sty`.
\makeatletter
\newfontfamily\germainGreekfont{cmun}[
  Extension=.otf, NFSSFamily=germaingreek,
  UprightFont=*rm, ItalicFont=*ti, SlantedFont=*sl,
  BoldFont=*bx, BoldItalicFont=*bi, BoldSlantedFont=*bl]
\newXeTeXintercharclass\germain@greek
\def\germain@greekrange#1#2{%
  \@tempcnta=#1\relax
  \loop
    \XeTeXcharclass\@tempcnta=\germain@greek
    \ifnum\@tempcnta<#2\advance\@tempcnta\@ne\repeat}
\germain@greekrange{"0370}{"03FF}
\germain@greekrange{"1F00}{"1FFF}
\def\germain@greekprev{\rmdefault}
\def\germain@greekin{%
  \edef\germain@greekprev{\f@family}\fontfamily{germaingreek}\selectfont}
\def\germain@greekout{\fontfamily{\germain@greekprev}\selectfont}
\def\germain@greekjoin{%
  \XeTeXinterchartokenstate=\@ne
  \@tempcnta=\z@
  \loop
    \ifnum\@tempcnta=\germain@greek\else
      \XeTeXinterchartoks\germain@greek\@tempcnta=\expandafter{%
        \expandafter\germain@greekout\the\XeTeXinterchartoks\z@\@tempcnta}%
      \XeTeXinterchartoks\@tempcnta\germain@greek=\expandafter{%
        \the\XeTeXinterchartoks\@tempcnta\z@\germain@greekin}%
    \fi
    \ifnum\@tempcnta<4095\advance\@tempcnta\@ne\repeat}
% babel-french sets its own classes, and saves and restores the interchar
% state, each time a language is selected: join again after it does.
\addto\extrasfrench{\germain@greekjoin}
\addto\extrasenglish{\germain@greekjoin}
\AtBeginDocument{\germain@greekjoin}
\makeatother

\usepackage{geometry}
\geometry{a4paper,margin=25mm,marginparwidth=28mm}
\usepackage{marginnote}
\usepackage{xcolor}
\usepackage[normalem]{ulem}
\setlength{\emergencystretch}{3em}
\usepackage{hyperref}
\hypersetup{colorlinks=true,linkcolor=black,urlcolor=black,pdfborder={0 0 0}}

% --- Batch metadata ---------------------------------------------------------
% Taken as they stand by the rendering script, which shows them at the head of
% the reading view. They fix what the file claims to cover. The macro names are
% the parent project's, so that the scripts stay shared; \folder here means the
% volume's slug — `fr-9115` — and \pages counts Gallica views.

\newcommand{\folder}[1]{\gdef\grothFolder{#1}}
\newcommand{\batch}[1]{\gdef\grothBatch{#1}}
\newcommand{\pages}[2]{\gdef\grothFirst{#1}\gdef\grothLast{#2}}
\newcommand{\foldertitle}[1]{\gdef\grothFoldertitle{#1}}

% The holder's own shelfmark, printed as they write it: « Français 9115 ».
\newcommand{\shelfmark}[1]{\gdef\grothShelfmark{#1}}

% The dating, copied verbatim from the holder's catalogue — never inferred
% here. « XIXe siècle » is the BnF's; a pass that dates a leaf from its content
% says so in a \note{}, as its own inference.
\newcommand{\dating}[1]{\gdef\grothDating{#1}}

% The status notice, printed in the title block and repeated as a diagonal
% watermark on every page of the PDF. The manuscripts are in the public
% domain; what this line declares is not a right but a quality — an
% unverified first pass by a machine. The skills write the line; a file
% without it gets no watermark, so leaving it out is visible in the output.
\newcommand{\watermark}[1]{\gdef\grothWatermark{#1}}

\gdef\grothFolder{?}\gdef\grothBatch{}\gdef\grothFirst{?}\gdef\grothLast{?}%
\gdef\grothFoldertitle{}\gdef\grothShelfmark{}\gdef\grothDating{}\gdef\grothWatermark{}

% --- The résumé -------------------------------------------------------------
\newenvironment{resume}
  {\small\setlength{\parskip}{0.45em}}
  {\par\medskip\hrule\medskip}

% --- Keywords ---------------------------------------------------------------
% In English, closing the modernised reading's résumé: the modern vocabulary
% under which to search for what the leaves construct. The manifest extracts
% them to tag the volume — this line is the single source of the tags.
\newcommand{\keywords}[1]{\par\smallskip\noindent
  {\footnotesize\textsc{Keywords} --- \textit{#1}}\par}

% --- The critical apparatus -------------------------------------------------

% The Gallica view number, in the margin. This is the anchor that lets a
% disagreement over a formula be settled by looking at *one* image rather than
% a volume of 750. The site uses it to turn the facsimile as the transcription
% is scrolled. A view is one scanned image: usually one side of a leaf.
\newcommand{\page}[1]{%
  \marginnote{\footnotesize\color{gray!70}\textsf{v.\,#1}}%
  \label{page:#1}\ignorespaces}

% The BnF's pencilled foliation, where the leaf carries it and a pass can read
% it: « 198r ». This is what the literature cites — « ff. 198r–208v » — and
% the one line that turns a citation into a view. Recorded, never inferred:
% a folio a pass cannot read is not written.
\newcommand{\folio}[1]{%
  \marginnote{\footnotesize\color{gray!50}\textsf{f.\,#1}}\ignorespaces}

% The modernised reading's coarser counterpart to \page{}: which views a
% section is drawn from.
\newcommand{\pagerange}[2]{%
  \marginnote{\footnotesize\color{gray!70}\textsf{v.\,#1--#2}}%
  \ignorespaces}

% Illegible: never guessed.
\newcommand{\ill}{\textcolor{red!60!black}{[\,\ldots\,]}}

% A doubtful reading: the word is given, and flagged as uncertain.
\newcommand{\uncertain}[1]{\uline{#1}}

% An editorial addition — a missing letter, an implied word.
\newcommand{\add}[1]{\textcolor{blue!50!black}{[#1]}}

% Struck out by the writer. What was crossed out is part of the document.
\newcommand{\struck}[1]{\sout{#1}}

% The transcriber's note, on the state of the page or the mathematics.
\newcommand{\note}[1]{\par\smallskip{\small\color{gray!60!black}\textit{[#1]}}\par\smallskip}

% A marginal note of hers, distinct from the body of the text.
\newcommand{\marginal}[1]{\par\smallskip\noindent\hspace{1em}%
  {\small\color{gray!60!black}#1}\par\smallskip}

% --- Title ------------------------------------------------------------------
% The title block is French, like the documents themselves.

\AddToHook{shipout/background}{%
  \ifx\grothWatermark\empty\else
    \begin{tikzpicture}[remember picture, overlay]
      \node[rotate=52, scale=3.4, align=center, text=gray!18,
            font=\sffamily\bfseries]
        at (current page.center) {\grothWatermark};
    \end{tikzpicture}%
  \fi}

\AtBeginDocument{%
  \begin{center}
    {\large\bfseries Manuscrits de mathématiciens (Gallica) ---
      \ifx\grothShelfmark\empty\grothFolder\else\grothShelfmark\fi}\\[2pt]
    {\itshape\grothFoldertitle}\\[4pt]
    {\small vues \grothFirst--\grothLast
      \ifx\grothBatch\empty\else\ (lot \grothBatch)\fi}\\[2pt]
    \ifx\grothDating\empty\else
      {\small\color{gray!60!black}Datation du catalogue : \grothDating}\\[2pt]
    \fi
    \ifx\grothWatermark\empty\else
      {\small\bfseries\color{red!45!gray}\grothWatermark}\\[2pt]
    \fi
    {\footnotesize\color{gray!60!black}Transcription établie d'après les images IIIF de Gallica.
     Source gallica.bnf.fr / Bibliothèque nationale de France.\\
     Les lectures incertaines sont soulignées, les illisibles notées [\,\ldots\,],
     les ajouts entre crochets ; \textsf{v.} numérote les vues de Gallica,
     \textsf{f.} la foliotation au crayon de la BnF.}
  \end{center}
  \vspace{1em}}

\endinput


\folder{naf-9544}
\shelfmark{NAF 9544}
\batch{1}
\pages{1}{20}
\dating{1601-1700}
\watermark{Lecture automatique\\non vérifiée}
\foldertitle{Recueil de lettres autographes de D'ALEMBERT, l'abbé BOSSUT, Sophie GERMAIN, C. HUYGHENS et le P. MERSENNE.}

\begin{document}

\page{3}
\note{Page de titre de l'album, d'une main calligraphique de la fin du XIX\textsuperscript{e} siècle ; la page de gauche, le contreplat, est blanche. En haut à droite, dans un cartouche : « Nouv. Acq. Fr. 9544 ». Au milieu : « Autographes de Savants », le dernier mot souligné. En bas : « Volume de 36 Feuillets — 27 Novembre 1899 ». Rien d'autre sur la page.}

\page{4}
\note{Premier feuillet de montage de l'album, sans texte manuscrit ancien. En haut à droite, à l'encre, le chiffre 1 (voir l'en-tête) ; au-dessous, la cote « N. a. franç. 9544. » soulignée d'une accolade, et au milieu un grand « 207 » ; en haut, vers le milieu, une marque brève (« P » ou « 1° »). Y est collée une coupure de catalogue de vente imprimée, qui annonce la pièce des vues 5–6 : « 3. Alembert (Jean Lerond d'), illustre savant et philosophe du XVIII\textsuperscript{e} siècle, de l'Acad. fr. — Let. aut. (minute) à Malesherbes ; (1765), 2 p. 1/4 in-8. — Il désire n'être plus pensionnaire surnuméraire dans la classe de géométrie, mais pensionnaire ordinaire dans la classe de mécanique ; c'est le vœu de l'Académie des Sciences ; il prie Malesherbes de parler en ce sens à M. de Saint-Florentin. » En marge de la coupure, à gauche, la marque manuscrite « R. D. 9555 » soulignée d'une accolade. La page de gauche, revers de la page de titre, est blanche.}

\page{5}
\note{Petit feuillet monté sur onglet, de la main de d'Alembert, écrit à la troisième personne : une minute ou un mémoire plutôt qu'une lettre (la coupure de catalogue de la vue 4 dit « Let. aut. (minute) à Malesherbes ; (1765) »). En haut à gauche, d'une main plus petite et plus tardive, « a M. de Malesherbes » ; au-dessus du titre, « 1765 », de la même petite main semble-t-il, le premier chiffre caché sous le paraphe de la capitale. En haut à droite, à l'encre, le chiffre 2 (voir l'en-tête). En bas à gauche, au crayon ou à l'encre, la marque « R. D. 9555 » soulignée d'une accolade, la même que sur la coupure de la vue 4. Le cachet rond de la Bibliothèque nationale est apposé au milieu du texte.}

\uncertain{Lettre de} M. d'alembert
\note{titre souligné. Les lettres qui suivent « Lett » se lisent « rannee » plutôt que « re de » ; la lecture est offerte par le sens.}

M. D'alembert est pensionnaire surnumeraire depuis 9 ans et demi dans la classe de Geometrie, il demande, \& l'academie demande pour lui, qu'il passe à la place de pensionnaire ordinaire mechanicien; la seule objection qu'on puisse lui faire, et qu'on lui fasse en effet, est celle de la difference des classes; objection sans fondement puisqu'il y a mille exemples dans l'academie, de passages d'une classe dans une autre, et surtout de la classe de Geometrie dans la classe de mechanique, \& reciproquement, ces deux classes ayant toujours roulé ensemble. M\textsuperscript{r}. de montigni en 1758 a passé de la place d'associé Geometre à celle de pensionnaire mechanicien; M\textsuperscript{r}. d'alembert demande beaucoup moins, puisqu'il est deja pensionnaire surnumeraire, \& qu'il ne s'agit que
\note{« Ordinaire » avec une capitale. La page s'arrête en pleine phrase ; la suite est au verso, vue 6.}

\page{6}
\note{Le petit feuillet de la vue 5 est ici déplié : à gauche son verso, qui continue la phrase ; à droite un second feuillet de même format, qui porte en haut à droite le chiffre 3 à l'encre et le cachet rond de la Bibliothèque nationale. Les deux pages forment un seul texte.}

de le faire changer de classe; ce qui ne prejudicie a personne 1°. parce qu'il est plus ancien que tous les associés. 2°. parce que M\textsuperscript{r}. de Courtivron le plus ancien après lui, demande la veterance, et a declaré par lettres qu'il ne vouloit point concourir avec M\textsuperscript{r}. d'alembert; ainsi M\textsuperscript{r}. de vaucanson qui se trouve \struck{\ill{}} \add{le second} associé mechanicien, deviendra le premier, c'est tout ce qu'il pouvoit esperer, en supposant que les choses se passassent à l'ordinaire, \& que M\textsuperscript{r}. de Courtivron n'eût point demandé la veterance.
\note{« le second » est écrit au-dessus de la ligne, au-dessus d'un mot noyé sous une rature épaisse et suivi d'un point, qui n'a pu être lu.}

ainsi M\textsuperscript{r}. de S\textsuperscript{t}. florentin n'a autre chose à repondre à la demande de l'academie que \emph{oui} ou \emph{non}; M\textsuperscript{r}. d'alembert espere que Monsieur de malesherbes voudra bien representer à M\textsuperscript{r}.
\note{« oui » et « non » sont soulignés. Sur « la » de « à la demande », une petite surcharge. La page s'arrête ici ; la phrase continue sur la page de droite.}

de S\textsuperscript{t}. florentin, combien le \emph{non} seroit injuste, malhonnête pour l'academie qui doit savoir ce qui lui convient le mieux, \& on ose dire, odieux après les travaux \& les sacrifices de M\textsuperscript{r}. d'alembert
\note{« non » souligné. Un trait horizontal, sous la dernière ligne, ferme le texte. Ni signature ni date : le reste de la page est blanc.}

\page{7}
\note{Pas de texte de lettre sur cette vue. À gauche, le verso blanc du feuillet 3 (un léger report d'encre en haut). À droite, un feuillet de montage portant à l'encre, en haut à droite, le chiffre 4 ; la lettre de Bossut y est montée, pliée, et l'on en voit la partie rabattue. Sur celle-ci, au crayon, un grand « 96 » entouré d'un ovale ; à l'encre, au-dessus, un « 6 » ; au bas, d'une main moderne, le nom « Bossut » ; sur une bande de papier collée au-dessous, un « 16 » à l'encre. Aucun de ces chiffres n'est une foliation de la bibliothèque (voir l'en-tête).}

\page{8}
\note{Première page d'une lettre de l'abbé Bossut (la signature est à la vue 10), montée sous le rabat de la vue 7, qui paraît encore en haut de la page. La date est en haut à droite ; au-dessous, un trait, un point et le chiffre 5 à l'encre (voir l'en-tête), et, au-dessus du millésime, au crayon semble-t-il, un petit « 40 ». Au milieu de la page, le cachet rond de la Bibliothèque nationale ; en bas à gauche, la marque « R. D. 9555 » soulignée d'une accolade. L'écriture du verso transparaît, à l'envers, dans le haut de la page.}

paris 12 Janv. 177\uncertain{4}
\note{le dernier chiffre du millésime est ouvert, en forme de u ou de h : il se lit 4 plutôt que 6, sans certitude.}

Reponse donnée
\note{mention de réception, au milieu de la page, sous la date, d'une autre main et d'une autre plume, sans doute celle du destinataire.}

Monsieur le Comte

Je vous supplie de me pardonner, Monsieur le Comte, si j'ai un peu differé de répondre à la lettre dont vous m'avés honoré en datte du 29 novembre dernier. Ce delai a été occasionné par des affaires que je ne detaillerai pas ici, crainte de vous ennuyer, et parceque d'ailleurs je veux \uncertain{tenir} mon pardon uniquement de votre bonté. l'honneur que votre celebre academie m'a fait
\note{la main de la lettre forme l'r final en petite boucle ouverte, qui se lit o ou w (« Monsieur », « honneur », « ennuyer »). La phrase se poursuit au verso, vue 9.}

\page{9}
\note{La lettre de Bossut dépliée : à gauche, le verso de son premier feuillet (la date et la mention « Reponse donnée » du recto y transparaissent à l'envers, dans le haut) ; à droite, le recto du second feuillet, qui porte en haut à droite un petit « 41 » et, au-dessous, un « 6 » plus grand, à l'encre (voir l'en-tête). Le rabat de la vue 7 reste visible au-dessus.}

de m'adopter, me penetre de la plus juste et de la plus vive reconnoissance. Il est bien glorieux pour moi, de voir mon nom associé à ceux de tant de grands hommes, qui decorent votre liste! quels remerciments ne vous dois je pas, Monsieur le Comte, de m'avoir attiré cette faveur? j'ose vous en demander encore une: celle de me mettre aux pieds de l'académie, et de luy temoigner la profonde vénération que j'ai pour tous ses membres, et l'envie extrême que j'ai de meriter leurs bontés. Mon ambition seroit comblée, si la fortune me presentoit l'occasion
\note{« m'adopter » est écrit vite, les lettres du milieu mal formées ; le sens et la phrase de la vue 8 (« l'honneur que votre celebre academie m'a fait ») le donnent. La page s'arrête ici et la phrase continue sur la page de droite.}

de réaliser ces sentiments.

Je me suis acquitté, Monsieur le Comte, de votre commission auprès de M. D'Alembert. il m'a chargé de vous faire un million de remerciments et de compliments. vous avés emporté l'estime et les regrets de tous ceux qui ont eu le bonheur de vivre ici avec vous; je n'en suis pas surpris, puisqu'aux connoissances les plus profondes, vous reunissés toutes les qualités sociales. C'est pour moi un vrai chagin de n'avoir fait,
\note{« chagin » : le mot est écrit sans r visible, pour « chagrin ». La phrase continue au verso, vue 10.}

\page{10}
\note{Verso du second feuillet de la lettre de Bossut, à gauche ; la page de droite, le revers du feuillet de montage, est blanche. Le texte du recto transparaît. Le cachet rond de la Bibliothèque nationale est apposé à gauche de la signature.}

pour ainsi dire, que vous entrevoir. Et si jamais quelque circonstance heureuse me rapproche de vous, je me propose bien de vous cultiver, et de mettre tous mes soins à vous convaincre de la haute estime et du profond respect avec lequel je suis,

Monsieur le Comte,

votre très humble et très obeissant serviteur

Bossut
\note{signature, avec un paraphe ; un trait horizontal la surmonte. Fin de la lettre, commencée à la vue 8. Le destinataire, « Monsieur le Comte », n'est pas nommé : la lettre dit seulement qu'une « celebre academie », dont il a transmis l'élection, vient d'associer Bossut, et que le comte a « emporté l'estime et les regrets » de ceux qui ont vécu « ici » avec lui, d'Alembert compris (vue 9).}

\page{11}
\note{Feuillet de montage, sans texte manuscrit ancien  ; la page de gauche est blanche. En haut à droite, à l'encre, le chiffre 7 (voir l'en-tête), et tout à fait dans l'angle un « 109 » souligné d'une courbe. En haut, vers le milieu, « 17\textsuperscript{a} » à l'encre, et au-dessous un grand « 95 » au crayon. Au crayon encore, au-dessus de la coupure, « \uncertain{cat. Dufour} 30 janv. 1889 ». La coupure de catalogue imprimée annonce la lettre des vues 14–16 : « 19. BOSSUT (Charles), célèbre géomètre, membre de l'Institut, n. à Tartaras (Loire), 1730, m. 1814. — L. a. s. à Antide Janvier ; Paris, 2 novembre an I\textsuperscript{er} (1792), 2 p. 1/2 in-4. — Intéressante lettre où il le remercie de l'envoi du portrait de Lalande. Il l'engage à demander un logement aux Menus. » En marge, « R. D. 9555 » souligné d'une accolade ; au-dessous de la coupure, un « 20 » à l'encre.}

\page{12}
\note{À gauche, une page blanche. À droite, un nouveau feuillet de montage, en partie couvert par une fiche manuscrite collée par le haut, qui recouvre la lettre de Bossut des vues 14–16. Sur la fiche, en haut à gauche, un grand « 1498 » (ou « 4498 ») au crayon ; au milieu, « 17\textsuperscript{b} » à l'encre ; en haut à droite, un « 8 » à l'encre (voir l'en-tête) et, dans l'angle, un petit « 113 ». Au crayon, sur la bande de papier qui dépasse sous la fiche : « Arrigoni, XIII\textsuperscript{e} catalogue \uncertain{1878} ».}

C. Bossut. Savant Geometre membre de l'academie n. 1730 m 1814
\note{fiche d'une main de collectionneur, non de Bossut, sur deux lignes, le nom souligné.}

\page{14}
\note{La fiche de la vue 12 est ici soulevée : dessous, montée sur le feuillet, la première page d'une lettre de Bossut. Tout en haut à droite, au-dessus de la date et d'une plume plus fine, un petit « 9 » (voir l'en-tête). Le cachet rond « BIBLIOTHÈQUE NATIONALE » est apposé au milieu de la page ; en bas à gauche, « R. D. 9555 ». La lettre est celle qu'annonce la coupure de la vue 11 ; l'adresse de la vue 16 la dit adressée à « Monsieur Janvier, horloger mecanicien ».}

Paris 2 9\textsuperscript{bre} l'an 1. de la Rep.
\note{« 9\textsuperscript{bre} » : novembre.}

Je suis infiniment reconnoissant, Monsieur, de la bonté que vous avés eu de m'envoyer le portrait de M\textsuperscript{r}. de la Lande, qui me paroit très bien gravé et très ressemblant.
\note{« bonté » : l'initiale est repassée à l'encre plus noire.}

Il m'est impossible de vous donner aucun conseil par rapport à la petition que vous voulés presenter à l'assemblée nationale. Je vous avois même prié, et vous m'aviés promis, de ne pas m'envoyer la copie ou la minute de cette petition. Mais vous savés d'ailleurs, Monsieur, que j'estime infiniment vos talents; je crois que vous
\note{« talents » est surchargé, un t écrit par-dessus le début du mot. La phrase continue à la vue suivante.}

\page{15}
\note{La lettre à Antide Janvier dépliée : à gauche, le verso du premier feuillet ; à droite, le recto du second, qui porte en haut à droite, à l'encre, « 10 » et, au-dessous, un petit « 111 » (voir l'en-tête). La fiche de la vue 12 reste visible au-dessus. Une tache ronde dans la marge droite de la page de droite recouvre la fin de deux lignes.}

merités un logement dans les maisons nationales. Je vous engage à demander au ministre le logement que M\textsuperscript{r}. \uncertain{Paris} occupoit aux menus, en le priant de se faire rendre compte de vos talents par la commission du museum. alors, me trouvant à portée legalement de vous rendre service, vous verrés, Monsieur, le zele que j'y mettrai. si votre affaire ne prend pas cette marche, j'aurai peut être de la peine à reussir pour un autre objet en votre
\note{le nom qui suit « M\textsuperscript{r}. » se lit « paris » ou « pavis » ; la lettre supérieure de « M\textsuperscript{r}. » est mal formée. « museum » est écrit avec un s long.}

faveur. Les logements aux menus ne regardant pas en general la commission du museum; mais i\add{l} pourroit se faire que M\textsuperscript{r}. Roland desirat savoir \uncertain{la} façon de penser de la commission sur votre compte.
\note{la fin de « il » et le mot avant « façon » sont pris dans la tache de la marge ; « la » est offert par le sens.}

Recevés les salutations de votre concitoyen

Bossut
\note{signature à la suite de la dernière ligne, avec un grand paraphe. Fin de la lettre commencée à la vue 14. Le cachet rond de la Bibliothèque nationale est apposé sous la signature ; le reste de la page est blanc.}

\page{16}
\note{Le second feuillet de la lettre à Janvier, rabattu et vu par son verso : le panneau d'adresse, écrit en travers de la page (la vue doit être tournée d'un quart de tour pour être lue), avec la déchirure du cachet dans la marge. Le texte de la page de droite de la vue 15 y transparaît. Au-dessus de l'adresse, un « 2 » isolé à l'encre, de la même orientation. La page de gauche, un feuillet de montage, est blanche.}

A Monsieur

Monsieur Janvier, horloger mecanicien. aux menus, rue poissonniere et bergere

A Paris

\page{17}
\note{La page de gauche est blanche. À droite, deux pièces montées l'une au-dessus de l'autre. En haut, un petit billet de la main de Bossut, qui porte en haut à droite, à l'encre, un petit « 114 » (ou « 11h ») et « 11 » (voir l'en-tête), le cachet rond de la Bibliothèque nationale et, en bas à gauche, « R. D. 9555 ». Au-dessous, une chemise d'une main de collectionneur italien, transcrite après le billet.}

J'ai reçu du Citoyen agasse la 65\textsuperscript{e}, 66\textsuperscript{e} et 67\textsuperscript{e} livraisons de l'encyclopedie methodique. A Paris ce 3 fructidor, an 10\textsuperscript{e}.

Bossut
\note{signature avec un grand paraphe, sur la dernière ligne. Un reçu, non une lettre : ni destinataire ni formule. Le chiffre du jour est un 3 bouclé qui pourrait être un 5.}

\note{Ce qui suit est la chemise, feuillet plus grand monté sous le billet. En haut à gauche, un grand « 97 » au crayon, entouré d'un arc ; à côté, à l'encre, « 18 » et, au-dessous, « 7 » ; en haut à droite, un petit « 64 » et, au-dessous, « 12 » souligné (voir l'en-tête). Trois notices, séparées par des traits, dans une écriture italienne soignée.}

abate Carlo Bossut Geometra nato a Tartaras nel 1730, morto nel 1814

marchese Simone de Laplace Geometra nato a Beaumont-en-auge nel 1749 morto nel 1827

J. B. Giuseppe Delambre astronomo nato ad amiens nel 1749, morto nel 1822
\note{« Tartaras » : l'initiale est tracée comme un Z ou un 2. Les deux notices qui suivent celle de Bossut annoncent les pièces des vues 18 et 19.}

\page{18}
\note{Une nouvelle pièce, montée seule sur la page de droite ; la page de gauche est blanche. C'est le sommaire, de la main de Bossut à ce qu'indique une note ancienne en tête, d'un mémoire sur l'équilibre des voûtes, avec deux attestations de lecture dans la marge gauche. En haut à droite, à l'encre, un petit « 65 » et, au-dessous, « 13 » (voir l'en-tête). Le cachet rond de la Bibliothèque nationale est au milieu de la page ; en bas à gauche, « R. D. 9555 ».}

\note{En tête, à gauche, d'une main plus tardive, une note de collectionneur : « écriture de M. l'abbé Bossut », le nom souligné.}

21 Messidor, an 9\textsuperscript{me} de la Rep.

Lû le 21 messidor an 9 à la réserve des calculs
\note{dans la marge gauche, en haut, sur un fragment de papier aux bords irréguliers, d'une autre main que le sommaire et d'une plume plus appuyée ; donné ici avant le texte, qu'il ne coupe pas. Au-dessous, au crayon, d'une main moderne : « aut. de Laplace ».}

Le 21 messidor an 9 le cit. Bossut a lû le memoire dont cette feuille est le sommaire

Delambre
\note{dans la marge gauche, plus bas, d'une troisième main, signée « Delambre » avec un paraphe.}

Le memoire du C. Bossut sur l'équilibre des voutes est divisé en trois sections, dont la premiere traite de l'équilibre des voutes en berceau, la seconde de l'équilibre des voutes en domes, la troisieme contient des experiences et des observations sur la resistance des pierres, et sur la maniere dont les \struck{\ill{}} \struck{\uncertain{voutes}} rompent les \struck{\ill{}} \add{voutes} trop chargées.
\note{à la fin de la deuxième ligne, un mot biffé puis noyé sous un gribouillis en boucles ; au début de la ligne suivante, un mot raturé en boucles, qui paraît être « voutes » ; « voutes » est ensuite récrit au-dessus d'un troisième mot biffé, entre « les » et « trop ». L'ajout interlinéaire est de la main du texte.}

Section 1. Trois chapitres.

Chap. I. Détermination de l'épaisseur des pieds droits, lorsque la voute \uncertain{se rompt en} des points donnés des reins.

Chap. II. Relations générales entre la figure de la voute et les forces qui pressent les voussoirs.

Chap. III. De la figure qu'il faut donner aux faces latérales et extérieures des pieds droits, lorsqu'ils peuvent se rompre par tranches horizontales.
\note{la page s'arrête ici ; le sommaire continue au verso, vue 19.}

\page{19}
\note{Verso du feuillet de la vue 18, rabattu à gauche : la suite du sommaire. La page de droite est blanche. Dans la marge gauche, à mi-hauteur, une tache d'encre en long trait, sans lettre lisible.}

Section II. Trois chapitres.

Chap. I. Détermination de l'épaisseur qu'il faut donner aux pieds droits des voutes en dome, pour resister à la poussée de la partie superieure formant une calotte que l'on considere comme un corps continu.

Chap. II. Conditions de l'équilibre entre toutes les parties d'une voute en domes.

Chap. III. De la figure qu'il faut donner aux faces latérales et exterieures des pieds droits des voutes en domes, lorsque ces pieds droits peuvent se diviser par tranches horizontales.
\note{« la figure » est souligné d'un trait épais, « la » repassé.}

Section III. en deux chapitres.

L'un sur les experiences, l'autre sur les observations.
\note{fin du sommaire ; le reste de la page est blanc, sans signature. Le sommaire annonce les trois sections du mémoire mais ne donne pas les calculs que réserve l'attestation de la vue 18.}

\page{20}
\note{La page de gauche est blanche. À droite, les pièces qui annoncent l'autographe de Sophie Germain, lequel commence à la vue suivante. En haut, une chemise de papier fort : en haut à gauche, « CLXXI » et « 2810 » à l'encre ; en haut à droite, « 65 » sur « 1 », séparés par un trait, et à côté « 14 » (voir l'en-tête) ; au milieu, un très grand « 677 » à l'encre ; au-dessous, d'une petite main, « Sofia Germain (autogr.) ». Au-dessous de la chemise est monté un billet d'une main de libraire ou de collectionneur, transcrit ici ; il porte en haut à droite « 65 » sur « 2 » et un « 15 » (ou « 13 ») au-dessus. Le cachet rond de la Bibliothèque nationale y est apposé.}

Sous ce pli l'autographe de Sophie Germain, acheté le lundi, 30 juin, (cabinet Dubrunfaut)

La couverture imprimée Charavay sera envoyée avec une brochure quelconque.

A. M.
\note{un trait ondulé sépare les deux parties du billet. « A. M. » : des initiales, avec un point, sans autre indication. L'année de l'achat n'est pas écrite.}

\end{document}
